% !TeX root = ../main.tex
\section{Đối tượng khách hàng}

\subsection{Người mới chơi eSports và sinh viên}

Đây là nhóm người dùng có nhu cầu chơi game cao nhưng chưa có nhiều kinh nghiệm chọn gear hoặc có ngân sách hạn chế. Họ dễ bị ảnh hưởng bởi quảng cáo, review trên mạng hoặc lời khuyên từ cộng đồng, từ đó mua sản phẩm không thật sự phù hợp với nhu cầu cá nhân.

Với nhóm này, Mutux giúp giảm rủi ro mua nhầm bằng cách cho phép thuê hoặc thử thiết bị trong một khoảng thời gian ngắn trước khi quyết định mua. Điều này đặc biệt phù hợp với sinh viên, vì họ thường không muốn bỏ ra 1--3 triệu đồng cho một món gear khi chưa chắc sản phẩm đó đáp ứng đúng nhu cầu.

\subsection{Game thủ, đội tuyển, câu lạc bộ và phòng game}

Nhóm này có nhu cầu tối ưu thiết bị để cải thiện trải nghiệm hoặc hiệu suất thi đấu. Họ quan tâm nhiều hơn đến độ phù hợp, độ ổn định và cảm giác sử dụng thực tế trong quá trình luyện tập hoặc thi đấu. Việc được thử nhiều mẫu gear khác nhau giúp họ đưa ra lựa chọn chính xác hơn trước khi mua.

Đối với các câu lạc bộ eSports, đội tuyển hoặc phòng game, Mutux có thể hỗ trợ thuê thử nhiều thiết bị cho thành viên trước khi mua số lượng lớn. Phòng game cũng có thể dùng nền tảng để thử nghiệm các dòng gear mới trước khi trang bị cho hệ thống máy.

\subsection{Shop gaming gear}

Các cửa hàng bán thiết bị gaming có thể đưa sản phẩm lên nền tảng để người dùng thuê thử. Đây là một kênh marketing và bán hàng mới, giúp shop tăng độ tiếp cận, tạo thêm doanh thu từ thiết bị trưng bày hoặc thiết bị nhàn rỗi, đồng thời tăng khả năng chuyển đổi từ thuê thử sang mua thật.
